% Generated draft: requires mathematical and source review.
\documentclass[11pt]{article}
\usepackage[margin=25mm]{geometry}
\usepackage{fontspec}
\setmainfont{Latin Modern Roman}
\usepackage{amsmath,amssymb,mathtools,unicode-math}
\setmathfont{Latin Modern Math}
\usepackage{enumitem,needspace,xurl,hyperref}
\setlength{\parindent}{0pt}
\setlength{\parskip}{5pt}
\setlength{\emergencystretch}{3em}
\newcommand{\N}{\mathbb{N}}
\newcommand{\Z}{\mathbb{Z}}
\newcommand{\Q}{\mathbb{Q}}
\newcommand{\R}{\mathbb{R}}
\newcommand{\C}{\mathbb{C}}
\newcommand{\F}{\mathbb{F}}
\newcommand{\A}{\mathbb{A}}
\newcommand{\PP}{\mathbb{P}}
\newcommand{\E}{\mathbb{E}}
\newcommand{\eps}{\varepsilon}
\newcommand{\dd}{\,\mathrm{d}}
\begin{document}
\section*{UnsolvedMath ID: 800143}
\textbf{Source problem code:} catalogue entry
\subsection{Definitions and mathematical statement}
For a positive integer \(n=\prod_p p^{v_p(n)}\), define
\[
\Omega(n)=\sum_p v_p(n),\qquad \Omega(1)=0.
\]
For real \(x\geq 1\), put
\[
S(x)=\sum_{\substack{n\in\Z\,;\ 1\leq n\leq x}}(-2)^{\Omega(n)},
\qquad
\alpha=\limsup_{x\to\infty}\frac{|S(x)|}{x}.
\]
Here the sum is over positive integers, and the limsup is initially understood in the extended nonnegative real numbers. These domain conventions make explicit the standard interpretation of the supplied notation.

\textbf{Question attributed to Daniel R. Johnston.} Determine the exact value of \(\alpha\).

The supplied source reports that an unnamed, then forthcoming preprint by students at UNSW--Canberra establishes that \(\alpha\) exists and gives a large upper bound. Since an extended limsup is already defined, the intended additional assertion appears to be finiteness; neither the preprint nor its bound is supplied. This report does not assert convergence of \(|S(x)|/x\). At \(x=2^k\), for positive integers \(k\), the jump of \(S\) is \((-2)^k\), with absolute magnitude \(2^k\). The source states that these jumps imply
\[
\alpha\geq \frac12.
\]
It also attributes to reference label \texttt{sun2016pair}, Conjecture 1.1, the conjecture
\[
\alpha\leq 1,
\]
and reports extensive computational support under reference label \texttt{mossinghoff2021oscillations}. These are supplied source claims, not independently verified conclusions.

\textbf{Additional, incomplete contribution attributed to Greg Martin.} The record also contains a separate discussion of sign changes of error terms such as \(\psi(x)-x\). It describes repeated averaging of an explicit formula to make its largest term dominate the remaining terms, thereby obtaining lower bounds for sign changes from that term. It attributes to Kaczorowski, reference label \texttt{K1}, the unconditional bound
\[
\liminf \frac{W^\psi(T)}{\log T}\geq \frac{\gamma_1}{\pi},
\]
and to Morrill, Platt, and Trudgian, reference label \texttt{MPT}, improving work of Kaczorowski under reference label \texttt{K2}, the bound
\[
\liminf \frac{W^\psi(T)}{\log T}\geq \frac{\gamma_1}{\pi}+1.867\cdot 10^{-30}.
\]
The supplied excerpt does not define \(\psi\), \(W^\psi\), or \(\gamma_1\), does not explicitly specify the limiting regime in these displays, and does not state a question for this contribution. Thus these displays are retained as source context rather than as a fully specified second problem. The contributor describes the latter bound as the best improvement known to him at the time of the source. The source further says that analogous arguments extend to many other error terms, generally requiring the relevant \(L\)-functions to have no real zeros, and that GRH tends to be necessary for biased error terms. These are qualified contextual remarks, not additional assumptions on the question about \(\alpha\).
\subsection{Short English statement}
Determine the exact limsup of \(|S(x)|/x\), where \(S(x)\) sums \((-2)^{\Omega(n)}\) over positive integers \(n\leq x\), and \(\Omega(n)\) counts prime factors with multiplicity. The supplied source reports a lower bound of \(1/2\), a finite upper bound in an unidentified forthcoming preprint, and the conjecture \(\alpha\leq 1\). A separate appended discussion of sign changes lacks definitions and an explicit question.
\subsection{Sources}
\begin{itemize}
\item Supplied local UnsolvedMath registry record: permanent ID \(800143\), exact source code \texttt{catalogue entry}, title \emph{Comparative Prime Number Theory --- Problem 12}. The registry labels the record \texttt{open}; this is a supplied status, not an independently verified mathematical assessment.
\item Hamieh, Kadiri, Martin, and Ng, \emph{Comparative Prime Number Theory Problem List} (2024), source item Problem (\#12), \url{https://arxiv.org/abs/2407.03530}. The registry identifies these authors in its proposer field. Its provenance records extraction by \texttt{arxiv\_latex} and an access date of 2026-10-16. That access date is supplied provenance; no independent source visit or verification is claimed here.
\item The question about \(\alpha\) is attributed within the supplied statement to Daniel R. Johnston, University of New South Wales--Canberra at ADFA. The reported existence and upper-bound result is attributed to unnamed students at UNSW--Canberra in a preprint described at the time as soon to be released; no identifying citation is supplied.
\item Reference label \texttt{sun2016pair}, Conjecture 1.1: cited by the supplied statement for the conjecture \(\alpha\leq 1\). Full bibliographic information is absent.
\item Reference label \texttt{mossinghoff2021oscillations}: cited by the supplied statement for extensive computational support. Full bibliographic information is absent.
\item The appended sign-change discussion is attributed to Greg Martin, University of British Columbia. It cites Kaczorowski under reference labels \texttt{K1} and \texttt{K2}, and Morrill, Platt, and Trudgian under reference label \texttt{MPT}. Full bibliographic information for these references is absent. The description of the best improvement is the contributor's source-time assessment.
\end{itemize}
\end{document}
